\section{Related Work}\label{sec:related_work}

Industry focus has predominantly been on approximate neighbor search systems, with FAISS ~\cite{FAISS_paper}, ScaNN ~\cite{scann_google_paper}, SONG ~\cite{song_paper}, RAFT~\cite{RAFT_paper} among notable examples. These support algorithms like HNSW ~\cite{hnsw_paper}, IVFPQ ~\cite{pq_paper}, CAGRA \cite{ootomo2023cagra} on CPU and GPU platforms. Termed model-free, these methods use unsupervised algorithms for partitioning space using existing item embeddings, offering flexibility for any item set. In contrast, our approach employs deep neural networks for a model-based search index, fully operational on GPUs. We integrate item indexes with neural network weights within a PyTorch or TensorFlow model, training during index construction and using the model for retrieval.

Recently with more performance and memory available on GPUs several publications have appeared considering model based nearest neighbor search such as \cite{MoL_paper, tay2022transformer, rajput2023recommender, zhang2023modelenhanced}. Mixture of logits (MoL) \cite{MoL_paper} in its production deployed form implements weighted combination of cosine similarities with neural network gates used to infer per distance component weights. The MoL paper does not provide information on examples of implementation of logits components, and which embedings have been used in production. We extend on top of MoL and introduce practical algorithms on how to learn components of MoL. 
Conversely, research by \cite{tay2022transformer, rajput2023recommender, zhang2023modelenhanced} has explored using transformers and generative techniques for search indexes. Unlike our system, which stores item embeddings directly, these studies create semantic structures through clustering and transformers to generate document IDs.

A lot of research has been focused on representation learning with works representing posts and users in social networks \cite{Que2Search_paper, nxtpost_paper, PinnerSage_paper}. As one of the components in MoL we have used approaches similar to \cite{Que2Search_paper, nxtpost_paper, PinnerSage_paper}, and additionally extended it with approaches for cold start infrequent users using clustering representations.
Several previous works have explored the concept of model live-updates, which we expand on in this paper. These works include Monolith~\cite{liu2022monolith}, PERSIA~\cite{persia_paper}, and XDL~\cite{XDL_paper}. In contrast, traditional search engines, as seen in Facebook Search EBR ~\cite{fb_search_embeddings_paper, Que2Search_paper} via Unicorn~\cite{unicorn_paper}, and Lucene~\cite{Lucene_chen2023endtoend}, have primarily focused on live-update functionality for unsupervised indexing techniques such as ~\cite{FAISS_paper}. To the best of our knowledge, our paper represents one of the pioneering efforts in the realm of retrieval-based techniques for live-updating TensorFlow (TF) or PyTorch model-based retrieval indexes at a large-scale production level, with high QPS demands.
%Most of the attention across the industry has been given to approximate neighbor search systems with some of the examples including FAISS ~\cite{FAISS_paper}, ScaNN ~\cite{scann_google_paper}, SONG ~\cite{song_paper}, RAFT~\cite{RAFT_paper}. Those frameworks support a range of approximate neighbor search algorithms such as HNSW ~\cite{hnsw_paper}, IVFPQ ~\cite{pq_paper}, CAGRA \cite{ootomo2023cagra} on CPU and GPU based approximate nearest neighbor search implementations. We may consider such approaches as model free approaches as those methods primarily look into unsupervised algorithms for space partitioning given existing item embeddings. The benefit of model free approaches is their flexibility and ease to readily apply to any set of items with embeddings. Our work looks into using deep neural networks to deploy as model-based search index containing all items fully on the GPU. We store items index along side of the neural network weights within a PyTorch or TensorFlow model. We train the model weights during index build and serve the model as retrieval index.
% \cite{tay2022transformer} used hierarchical k-means with LLM-based document representation and transformer encoders. \cite{rajput2023recommender} and \cite{zhang2023modelenhanced} developed transformer-based clustering indexes.

%On the other side authors of \cite{tay2022transformer, rajput2023recommender, zhang2023modelenhanced} looked into clustering with transformers and generative retrieval to represent search indexes. Instead of storing item embeddings as we do in {\systemname}, authors of \cite{tay2022transformer, rajput2023recommender, zhang2023modelenhanced} built a semantic structures using clustering and transformers to infer document IDs in generative manner. \cite{tay2022transformer} used hierarchical kmeans with underlying LLM based document representation and transformers encoders to represent model based search index. Later \cite{rajput2023recommender} and \cite{zhang2023modelenhanced} looked into learnt transformer-based clustering index. In our work we focused on exhaustive model based solution where model scans over every item in its index to firstly evaluate attribute based constraints and then perform exhaustive KNN, which to the best of our knowledge were not addressed in prior works together with model based indexes. Because we store item features as part of the index we are able to do a lightweight ranking of the retrieved items in place effectively merging retrieval and ranking stack into one single call to the model. For one of the representations we applied kmeans based learning solution for representation learning needs, while \cite{tay2022transformer, rajput2023recommender, zhang2023modelenhanced} primarily used clustering to encode document IDs.


